\section{Introduction}\label{sec:intro}

% G = General (current situation, with public evidence)
Fresh vegetables are among the most volatile items in Taiwanese household
spending. Typhoons and heavy rain regularly damage field crops and disrupt
shipments. An analysis by Taiwan's agricultural authority of six major
disasters in 2004--2009 found that average wholesale vegetable prices rose by
about 10--40\% in the first week after a disaster, 30--90\% in the second
week, and up to about 100\% in the third week~\cite{coa2010typhoon}. The
pattern persists: after Typhoon Gaemi in 2024, the average wholesale vegetable price
in Taipei reached NT\$57.5/kg, a five-year high~\cite{taiwannews2024gaemi}, and
in December 2024 vegetable prices in the consumer price index were 30.69\%
higher than a year earlier~\cite{ltn2025cpi}. During Typhoon Podul in August
2025, the wholesale price of green onions rose by about one third, and a stall
in a Taipei traditional market sold a single bundle for
NT\$380~\cite{ltn2025onion,pts2025onion}; that month, vegetables were 10.9\%
more expensive than a year before~\cite{taipeitimes2025cpi}. Price information
itself is plentiful: the Ministry of Agriculture publishes daily wholesale
quotations~\cite{moa2026trans}, and Taichung City publishes daily prices from
its 13 public retail markets~\cite{taichung2026retail}. Yet a shopper in front
of a stall sees only one number---today's price in this store---and must decide
whether to buy now or wait.

% H = However (challenges)
However, three properties of the problem make this decision hard to support
with data. First, public averages are not store prices: price changes pass from
wholesale to retail incompletely and
asymmetrically~\cite{ward1982asymmetry,meyer2004asymmetric}. Second, retail
prices are sticky: most prices stay unchanged for weeks or months and then
jump~\cite{bils2004some,nakamura2008five}. In our data, 86.7--92.8\% of
seven-day price pairs have exactly the same price at both ends, so a forecaster
that always predicts ``no change'' is hard to beat. Third, large language
models (LLMs) are attractive as a natural-language front end, but they can
generate fluent yet unsupported content~\cite{huang2025survey}, which is
unacceptable when the output is a price.

% I = In literature (categories of solutions, old to new)
In the literature, agricultural price forecasting has moved from statistical
models and early forecasting services built on open market
data~\cite{peng2015developing}, to tree
ensembles~\cite{breiman2001random,ke2017lightgbm}, deep sequence
models~\cite{zhang2024short}, and, most recently, LLM-based
forecasters~\cite{zhao2025vegetable}. In parallel, LLM agents have learned to
call external tools~\cite{yao2023react,schick2023toolformer}, and the Model
Context Protocol (MCP) now standardizes how agents connect to tools and
data~\cite{anthropic2024mcp}.

% J = Judgement (strengths and weaknesses)
These lines of work have complementary weaknesses. Forecasting studies predict
wholesale or market averages and often compare learned models with each other,
although large competitions show that simple methods are hard to
beat~\cite{makridakis2020m4} and that added complexity, including LLM
components, does not always help~\cite{zeng2023transformers,tan2024language}.
Agent studies, conversely, focus on language tasks and rarely address numeric
forecasting, where a model that computes or guesses numbers itself undermines
trust.

% K = Keypoint (design spirit, RQs, contributions)
This paper proposes the Fresh Market Price Forecast Agent (FMPA). Its design
spirit is: \emph{anchor on what the user sees, compute only with tools, and
evaluate against the strongest simple baseline.} We study three research
questions:
\begin{itemize}
\item \textbf{RQ1:} How well can the relative price change be forecast at
fixed horizons of 7 and 14 calendar days, compared with no-change and momentum
baselines?
\item \textbf{RQ2:} Does a weekly-momentum adjustment (mode B) add value over
pure store-price anchoring (mode A)?
\item \textbf{RQ3:} Is a tool-based LLM agent a feasible, low-effort interface
for consumer price queries?
\end{itemize}
The main contributions of this work are as follows:
\begin{itemize}
\item \textbf{C1 (RQ3):} a domain AI agent architecture in which the LLM only
parses intent, selects MCP tools, and explains results, while every price is
retrieved or computed by deterministic tools, with explicit errors when data
are missing;
\item \textbf{C2 (RQ1, RQ2):} a store-price-anchored forecasting formulation
with two modes, evaluated at fixed 7- and 14-day calendar horizons under a
leakage-safe protocol against strong naive baselines;
\item \textbf{C3 (RQ1):} an empirical finding on 13 public markets in Taichung
City that RandomForest and HistGradientBoosting models do not beat these
baselines, with a domain explanation based on price stickiness.
\end{itemize}

% L = List the organization
The rest of this paper is organized as follows. Section~\ref{sec:related}
reviews AI methods and domain knowledge and identifies the research gap.
Section~\ref{sec:method} presents the proposed scheme.
Section~\ref{sec:experiments} describes the experimental evaluation, and
Section~\ref{sec:conclusion} concludes the paper.
